% !TEX root = template.tex
\section{Processing Pipeline}

\begin{figure}[htbp]

  \centering
\includegraphics[width=\linewidth]{processing_pipeline.png}
\caption{Processing pipeline.}
\label{fig:pipeline}
\end{figure}

In this section, we're going to explain the structure of the pipeline, designed for the analysis and processing of the input doppler traces; a sketch of the pipeline is depicted in the flow diagram in Figure 1. In the following, all the processing steps:
\begin{enumerate}
\item \textbf{Input data:} The raw Doppler traces, associated with specific activities (defined by the labels E, L, W, R, J), serve as the input. These traces are acquired using the techniques proposed by [4] and [5], as detailed in Section IV;
\item \textbf{Phase sanitization and computing the Doppler Traces:} This step, using the techniques specified in Section IV, filters environmental and hardware noise from the raw data to ensure its reliability for subsequent analysis. After that, the input data is computed and the data sequence is obtained, as described in Section IV;
\item \textbf{Dataset transformation and creation:} Following the cleaning process, transformations are applied to the data to enhance the robustness of the model against variations in the input. The specific techniques applied in this step are detailed in Section IV;
\item \textbf{Deep neural architecture and Classification:} The transformed data are processed by a Resnet (whose architecture is defined in Section V) to predict the most probable class associated with the input. The system then outputs a loss value (logits) based on the predict.
\end{enumerate}

\section{Signals and Features}

\begin{figure}[htbp]
\centering
\includegraphics[width=\linewidth]{labels.png}
\caption{Spectograms of the five different activities and their associated labels, as proposed by [4]}
\label{fig:spectrograms}
\end{figure}

As described by [5] and [4], the dataset comprises WiFi Channel Frequency Response (CFR) data collected within an IEEE 802.11ac network. The signals were captured by a monitor node (an ASUS RT-AC86U router) while two terminals exchanged traffic on channel 42 operating at a center frequency of 5.21 GHz with an 80 MHz bandwidth while a human subject acted as a physical obstacle to the transmission by performing various activities. Due to constraints in computational resources, this experiment utilizes a subset of the labeled activities and subjects. Specifically, the evaluated activities are empty 'E', sitting 'L', walking 'W', running 'R', and jumping 'J' (Figure \ref{fig:spectrograms}).

The dataset treatment differs significantly depending on the classification task:
\begin{itemize}
    \item \textbf{Activity Recognition:} The dataset is split such that data from four subjects are allocated to the training and validation phases, while data from two completely distinct, unseen subjects (Subjects 6 and 7) are reserved exclusively for testing to evaluate true model generalization on unseen individuals.
    \item \textbf{Person Identification:} The problem is framed as a binary classification between two specific subjects (Subjects 4 and 5). The test set consists of disjoint tracking sessions from these same subjects to evaluate identification accuracy on new physical environment sessions rather than zero-shot unseen people. To robustly learn a highly discriminative latent space, the training dataset is dynamically wrapped into a Triplet configuration that generates Anchor, Positive, and Negative samples on the fly. Furthermore, a Weighted Random Sampler is applied during batch generation to perfectly balance the class distribution and stabilize contrastive learning.
\end{itemize}

About the computation of CFR, in [4], the authors mathematically formulate the raw CFR acquired from commercial Wi-Fi devices communicating via Orthogonal Frequency Division Multiplexing (OFDM). The multi-path CFR across subchannels is defined as:
\begin{equation}
H_k(n) = A_k(n)e^{j\phi_k(n)} = \sum_{p=0}^{P-1} A_p(n)e^{-j2\pi(f_c+k/T)\tau_p(n)}
\end{equation}
where $H_k(n)$ is the sub-channel response, $A_k(n)$ is attenuation, $\phi_k(n)$ is phase shift, $P$ is the path count, $A_p(n)$ is path attenuation, $f_c$ is carrier frequency, $T$ is OFDM symbol time, and $\tau_p(n)$ is path delay. This raw signal model contains non-ideal hardware phase artifacts that distort real-world measurements.

After that, there's a Sanitization phase, in which this raw CFR is processesed by removing the phase offset, outputting a cleaned signal $\hat{H}_k$, which is computed as:
\begin{equation}
\hat{H}_k \approx A_{p^*} e^{-j2\pi(f_c+k/T)\tau_{p^*}} H_k
\end{equation}
where $\hat{H}_k$ is the corrected CFR for sub-channel $k$, $A_{p^*}$ and $\tau_{p^*}$ are the amplitude and delay of the strongest reference path $p^*$, while $H_k$ represents the isolated dynamic multi-path component required for detecting human activity.

Finally the last step of the preprocessing phase is the Doppler Trace Computation, in which the sanitized CFR ($\hat{H}_k$) previously computed are transformed into Doppler spectrograms to isolate human movement velocity from static reflections. By arranging sanitized phase data into timesequenced observation matrices and applying a short-time Fourier transform (STFT), temporal variations are converted into a Doppler spectrogram. The physical movement model is defined as:
\begin{equation}
v_p \cos(\alpha_p) = \frac{u c}{f_c T_c N_D}
\end{equation}
where $v_p \cos(\alpha_p)$ is the directional velocity of scatterer $p$, $u$ is the discrete Doppler bin index associated with physical motion, $c$ is the speed of light, $f_c$ is the central carrier frequency, $T_c$ is the sampling period between packets and $N_D$ is the Number of zero-padded signal samples. Specifically, the computation of Doppler traces is performed with the following parameters:
\begin{itemize}
\item Number of channels estimates for Doppler computation per sample: 31
\item Number of packets for sliding operations: 1
\item Number of Doppler vectors per window: 340
\item Number of samples for window sliding: 30
\item Number of monitoring antenna: 4
\end{itemize}

Once the sanitized data are obtained, an \textbf{Instance Normalization} step is applied: every sample is independently normalized to zero mean and unit variance. This establishes crucial robustness against varying absolute signal reflections across different individuals (particularly vital for the unseen Subjects 6 and 7).

Following normalization, the following transformations are applied sequentially to the training dataset, in order to improve model generalization:
\begin{enumerate}
    \item \textbf{Gaussian noise:} Zero-mean additive white Gaussian noise ($\mu = 0$, $\sqrt{\sigma^2} = 0.1$) is injected into the data sequence to simulate signal fluctuations.
    \item \textbf{Uniform scale variation:} A scaling factor chosen uniformly from $[0.8, 1.2]$ is applied to randomly alter the apparent duration of the sequence.
    \item \textbf{Block masking:} A contiguous block representing 10\% (proportional ratio 0.1) of the total sequence length is randomly masked in both the time and frequency domains (similar to Cutout or SpecAugment) to train the model to handle missing data.
\end{enumerate}
To ensure unbiased evaluation, these augmentations are disabled during the validation and testing phases. Specifically, the Gaussian noise and block masking parameters are set to 0, and the scale variation factor is fixed to 1.0. These transformations are applied dynamically upon dataset instantiation directly in the GPU VRAM.
